\setchapterpreamble[u]{\margintoc}
\chapter{Некоторые понятия теории множеств}
\tikzsetfigurename{SetTheory_}

Ниже мы введём некоторые понятия из теории множеств, которые будут необходимы при изложении дальнейшего материала.
Мы будем полагать, что читатель знаком с понятием множества и лишь напомним некоторые моменты.
Более глубокое погружение в тему можно начать, например, с работы Николая Константиновича Верещагина и Александра Ханиевича Шеня <<Начала теории множеств>>~\sidecite{верещагин2000начала}.

\section{Основные определения}

Будем говорить\sidenote{Стоит отметить, что формальное определение множества --- не самая очевидная задача. Некоторые <<подводные камни>> обсуждаются, например, в уже упомянутой работе <<Начала теории множеств>> на страницах 29--31.}, что \emph{множество} $M$ состоит из \emph{элементов} $m$.
Факт, что элемент $m$ принадлежит множеству $M$ записывается как $m \in M$.
Отдельно введём обозначение $\varnothing$ для \emph{пустого множества} --- множества, не содержащего элементов.

\begin{definition}[Мощность множества]
    \emph{Мощность} множества $M$ --- это количество элементов в него входящих.
    Записывается как $|M| = k$: мощность множества\sidenote{Пустое множество не содержит элементов: $|\varnothing| = 0$} $M$ равна $k$.
\end{definition}

\begin{definition}[Подмножество]
    Множество $M_2$ является \emph{подмножеством} множества $M_1$ если все элементы $M_2$ являются и элементами $M_1$.
    Записывается как $M_2 \subseteq M_1$.
\end{definition}

\begin{definition}[Равенство множеств]
    Множество $M_2$ \emph{равно} множеству $M_1$ если они состоят из одних и тех же элементов ($M_2 \subseteq M_1$ и $M_1 \subseteq M_2$).
    Записывается как $M_2 = M_1$.
\end{definition}

\begin{definition}[Собственное подмножество]
    Множество $M_2$ является \emph{собственным подмножеством} множества $M_1$ если оно является подмножеством и при этом $M_2 \neq M_1$.
    Записывается как $M_2 \varsubsetneq M_1$.
\end{definition}

\begin{definition}[Булеан]
    \emph{Булеан} или \emph{множество всех подмножеств} множества $M$ --- это множество, состоящее из всех подмножеств $M$, включая пустое множество $\varnothing$ и само множество $M$.
    Обозначается как $2^M$.
\end{definition}

Напомним также определение основных операций над множествами (или же теоретико-множественные операций).

\begin{definition}[Объединение множеств]
    Говорим, что множество $S$ является объединением множеств $S_1$ и $S_2$, если
    \[S = \{ x \mid x \in S_1 \text{ или } x \in S_2\}.\]
    Записывается как $S = S_1 \cup S_2$.
\end{definition}

\begin{definition}[Пересечение множеств]
    Говорим, что множество $S$ является пересечением множеств $S_1$ и $S_2$, если
    \[S = \{ x \mid x \in S_1 \text{ и } x \in S_2\}.\]
    Записывается $S = S_1 \cap S_2$.
\end{definition}

\begin{definition}[Разность множеств]
    Говорим, что множество $S$ является \emph{разностью} множеств $S_1$ и $S_2$, если
    \[S = \{ x \mid x \in S_1 \text{ и } x \notin S_2\}.\]
    Записывается как $S = S_1 \setminus S_2$.
\end{definition}

\begin{definition}[Дополнение множества]
    Если множество $S_1$ является подмножеством $S$, то $\overline{S} = S \setminus S_1$ --- это \emph{дополнение} множества $S_1$.
\end{definition}

\begin{definition}[Симметрическая разность множеств]
    Говорим, что множество $S$ является \emph{симметрической разностью} множеств $S_1$ и $S_2$, если
    \[S = (S_1 \setminus S_2) \cup (S_2 \setminus S_1).\]
    Записывается как $S = S_1 \bigtriangleup S_2$.
\end{definition}

%\sidenote{
%        Обратите внимание, что в определении пары упорядочены.
%        То есть декартово произведение множеств --- не является коммутативной операцией.}

\begin{definition}[Декартово произведение множеств]
    Множество $S$ является декартовым произведением множеств\sidenote{ Обратите внимание, что в определении пары упорядочены. То есть декартово произведение множеств --- не является коммутативной операцией.} $S_1$ и $S_2$, если
    \[S = \{ (x,y) \mid x \in S_1, y \in S_2 \}.\]
    Записывается как $S = S_1 \times S_2$.
\end{definition}

Данное выше определение часто обобщается на произвольное количество множеств в том смысле, что декартово произведение $k$ множеств --- это упорядоченный кортеж\sidenote{
    Если действовать по данному выше определению, то получится, скорее, следующая конструкция. $A_0 \times A_1 \times \ldots \times A_{k-1} = (\ldots(A_0 \times A_1) \times \ldots \times A_{k-1}) = \{ (\ldots(a_0, a_1),a_2), \ldots, a_{k-1} \mid a_i \in A_i\}$.
    То есть мы получим какую-то конструкцию из пар, что не всегда удобно при дальнейшей работе.
    Правда, вопрос расстановки скобок (и, соответственно, того, какую именно конструкцию мы получим) --- тема для отдельного рассуждения.}
размера $k$.
Далее мы будем пользоваться таким обобщением.

Заметим, что приведённые выше операции обладают разными свойствами, на которые полезно обращать внимание при работе с ними.
Так, операция объединения является \emph{идемпотентной}, а разность не является \emph{коммутативной}.

Кроме чисто теоретико-множественных операций, нам понадобятся операции, которые что-то делают с элементами множеств.

\begin{definition}[Поэлеиентная операция над множествами]
    Пусть дано множество $S$ с определённой на нём операцией $\odot: S \times S \to S$, $S_1 \subseteq S$, $S_2 \subseteq S$, тогда
    \[S_1 \odot S_2 = \{ s_1 \odot s_2 \mid s_1 \in S_1, s_2 \in S_2\}.\]
\end{definition}

\begin{definition}[Степень множества]
    Пусть дано множество $S$ с определённой на нём операцией $\odot: S \times S \to S$, $S_1 \subseteq S$, тогда
    \[S_1^n = \{ \underbrace{s_1 \odot s_1 \odot \dots \odot s_1}_{\text{$n$ раз}} \mid s_1 \in S_1\}.\]
    При этом\sidenote{В данном случае нулевая степень даёт единицу, как мы и привыкли.} $S_1^0 = \{\varepsilon\}$.
\end{definition}

\begin{definition}[Замыкание множества]
    Пусть дано множество $S$ с определённой на нём операцией $\odot: S \times S \to S$, $S_1 \subseteq S$, тогда
    \[S_1^* = \bigcup_{n = 0}^{\infty} S_1^n.\]
\end{definition}


\section{Отношения}

\begin{definition}[Отношение]
    Любое подмножество $R$ декартова произведения $k$ множеств, для $k > 1, k < \infty$, называется \emph{$k$-местным отношением}.
\end{definition}

Если есть $k$-местное отношение $R$ и $(x_0,\ldots, x_{k-1}) \in R$, то часто говорят, что $x_i$ \emph{находятся в отношении} $R$.

Если $k=2$ и $A_0 = A_1 = A$, то говорят о emph{бинарных отношениях} над $A$.

\begin{definition}[Рефлексивное отношение]
    Бинарное отношение $R$ над $A$ является \emph{рефлексивным}, если для любого $a \in A$ $(a,a) \in R$.
\end{definition}


\begin{definition}[Транзитивное отношение]
    Бинарное отношение $R$ над $A$ является \emph{транизитивным}, если для любых $a_0 \in A, a_1 \in A, a_2 \in A$ , если $(a_0,a_1) \in R$ и  $(a_1,a_2) \in R$, то $(a_0,a_2) \in R$.
\end{definition}

\begin{definition}[Транзитивное замыкание отношения]
    Пусть $R$ --- бинарное отношение над $A$.
    Тогда транзитивным замыканием $R$ является наименьшее по мощности транзитивное отношение $R_t$ над $A$, такое, что\sidenote{Помним, что отношение --- это тоже множество.} $R \subseteq R_t$.
\end{definition}

\begin{definition}[Рефлексивно-транзитивное замыкание отношения]
    Пусть $R$ --- бинарное отношение над $A$.
    Тогда рефлексивно-транзитивным замыканием $R$ является наименьшее по мощности транзитивное и рефлексивное отношение $R_{rt}$ над $A$, такое, что $R \subseteq R_{rt}$.
\end{definition}
